\section{What this earns, prior mechanisms, and preserved boundaries}

The new operational consequence is a resource choice for an acquired joint
first-event law. A fitted-model route can learn numerical energy and
kinetic tables and use a marked-law stability theorem. The present route
instead keeps calling faithful local cells to implement the true kernels;
it need not learn a fine operating-rate table, and under the pairwise
coarse certificate need not enumerate that table's $2^d$ contexts at all.
The optional fitted-static hybrid learns only the invariant weights
finely and supplies true kinetics experimentally. These routes support
the same finite-menu bounded-loss endpoint with different acquisition,
per-record service and metrology costs. No one route is universally
cheaper, and none implies that a natural system's closed model has been
identified from short observations.

Bernoulli factories already transform samples from unknown probabilities
into samples of functions of those probabilities; analytic-factory
constructions and their finite tail theory are established
\cite{nacuperes}. Two-coin methods for intractable transition probabilities
and Bernoulli races are likewise established \cite{barker,bernoullirace}.
The wider idea that sampling can replace exact numerical evaluation is
also present in exact-approximate and pseudo-marginal MCMC
\cite{pseudo} and the exchange algorithm \cite{exchange}.
Thus a generic slogan that parameters need not be estimated before
sampling is not the contribution.

The specific construction here turns the \emph{permitted local physical
interventions} into the static and dynamic interfaces of a finite
rare-first-event operation. It closes four concrete interfaces:
\begin{enumerate}
\item A balanced artificial-field observation becomes an operating Gibbs
conditional, while paying the actual field shift and finite relaxation.
\item An actual operating first-exit race becomes a uniformized transition
coin, without equating biased and operating clocks or paying a rare full
dwell. Arbitrary locally requested state pins and timing are charged.
\item Unknown terminal weights in heavy-prefix counting are replaced by
a path of measured local log odds, so the rare static interface is earned.
\item Conditional query errors and conditional means are propagated
through static sampling, Palm calibration, marked rejection and decisions;
physical iid records are not assumed.
\end{enumerate}
The accepted cut and Palm operations supply the global rarity removal;
this construction does not claim to replace their strong static and
mixing premises with mere local observations.
No historical firstness is asserted for the combined theorem.

\subsection{Useful failed routes and why they are not used}

\begin{enumerate}
\item Coupling two full lifetime trajectories step by step accumulates
error over an enormous waiting time. It does not give the present
no-inverse-rarity operation. The completed finite block sampler skips
that lifetime, while the present procedure implements its short queries.
\item A bounded-time passive record cannot relatively estimate an
arbitrarily small local exit rate. The censored-exit factory does not
estimate that rate; it samples a coin whose correspondingly rare outcome
has the correct probability. The externally priced state pin is essential.
\item An artificial balancing field may change the total relaxation
clock. Using that observed clock as the operating clock would be wrong.
Every operating coin here uses the actual requested operating environment.
\item Coarse energy learning does not make field precision free.
Unbiasing with $e^{u-z}$ requires accurate knowledge and reproduction of
the actual additive shift. A hardware error floor yields an output
error floor unless another interface is proved.
\item Fast Gibbs sampling alone gives additive estimates of rare event
occupancy. Relative pin masses require the heavy-prefix identity and
the measured weight-ratio path. Counting a cut by waiting for it would
restore inverse rarity.
\item Reusing past random bits or observed cell states as though they
were independent fresh outputs is unjustified. Uniform conditional
kernel bounds, holds and conditional-mean concentration are the repairs.
\item Coupling every high-level calibration call at error $\delta/D$
can force a polynomial inverse-confidence static bill. The conditional
mean lemma retains a fixed per-proposal bias budget while sample counts
carry $\log(1/\delta)$. The former wrapper is a valid but weaker option.
\end{enumerate}

\subsection{Execution status}

The companion \texttt{factory\_audit.py} executes exact rational
finite-walk and two-coin-factory audits, including extremely small
operating-rate examples. It checks finite success probabilities against
the target identities and rigorous truncation bounds, and records the
finite service parameters in \texttt{factory\_results.json}.
The executed audit has 10 exact finite-walk cases, 20 exact Gibbs-factory
cases and 4 independently enumerated pin-mass identities. All checks
passed. It includes a positive operating rate $2^{-100}$ and, separately,
a zero-rate boundary diagnostic outside the positive-clock theorem.
These checks exercise the two local factories and their finite caps.
They do not execute the full interacting cut enumerator, the full physical
first-event controller, hardware calibration, or an actual experiment.
The theorem and its resource bill, rather than a small simulation, earn
the general consequence.

Every accepted original remains unchanged. The earlier acquired relative
CDF, independent first-event record sampler and finite-bit clock,
active acquisition/preparation/prevention, fixed-law prediction and
fixed-mechanics thermal operations remain at their original scopes.
The earlier all-strength quantum laws and thermal effusion result are not
reinterpreted by this classical construction.

\begin{thebibliography}{9}
\bibitem{accepted}
\raggedright
Astra research knowledge, \emph{First-event package and preserved sources},
revision \texttt{719e201332992002528c221c49f80f69431485f6}.
\url{https://github.com/harjassand/astra-research-knowledge/tree/719e201332992002528c221c49f80f69431485f6/research/returns/2026-11-05-p3-first-event}.
Exact local governing files used:
\texttt{next4\_root/first\_event\_report.tex},
\texttt{next4\_root/block\_completion.tex},
\texttt{next4\_root/finite\_compiler.tex},
\texttt{next4\_intervention/active\_acquisition\_certificate.tex},
and \texttt{next3\_root/correlated\_reliability\_report.tex}.
\bibitem{nacuperes}
S. Nacu and Y. Peres, \emph{Fast simulation of new coins from old},
Annals of Applied Probability 15(1A), 93--115 (2005).
doi:10.1214/105051604000000549.
\url{https://arxiv.org/abs/math/0309222}.
\bibitem{barker}
F. B. Gon\c{c}alves, K. \L atuszy\'nski and G. O. Roberts,
\emph{Barker's algorithm for Bayesian inference with intractable likelihoods},
Brazilian Journal of Probability and Statistics 31(4), 732--745 (2017).
doi:10.1214/17-BJPS374.
\url{https://arxiv.org/abs/1709.07710}.
\bibitem{bernoullirace}
S. Dughmi, J. D. Hartline, R. Kleinberg and R. Niazadeh,
\emph{Bernoulli Factories and Black-Box Reductions in Mechanism Design},
STOC 2017. doi:10.1145/3055399.3055492.
\url{https://www.cs.cornell.edu/~rdk/papers/stoc17-rad.pdf}.
\bibitem{pseudo}
C. Andrieu and G. O. Roberts,
\emph{The pseudo-marginal approach for efficient Monte Carlo computations},
Annals of Statistics 37(2), 697--725 (2009).
doi:10.1214/07-AOS574.
\url{https://arxiv.org/abs/0903.5480}.
\bibitem{exchange}
I. Murray, Z. Ghahramani and D. J. C. MacKay,
\emph{MCMC for doubly-intractable distributions}, UAI (2006).
\url{https://homepages.inf.ed.ac.uk/imurray2/pub/06doubly_intractable/}.
\end{thebibliography}
